\chapter{Calcolo degli Integrali Singoli}

Il calcolo integrale rappresenta l'operazione inversa della derivazione. In questo capitolo esploreremo le tecniche fondamentali per il calcolo degli integrali definiti e indefiniti.

\section{Integrali Notevoli}
Per procedere al calcolo di integrali complessi, è indispensabile conoscere le primitive delle funzioni elementari. Ecco un elenco delle più comuni:

\begin{itemize}
    \item \textbf{Potenza:} $\int x^n \, dx = \frac{x^{n+1}}{n+1} + C \quad (n \neq -1)$
    \item \textbf{Reciproco:} $\int \frac{f'(x)}{f(x)} \, dx = \ln|f(x)| + C$
    \item \textbf{Esponenziale:} $\int e^x \, dx = e^x + C$
    \item \textbf{Seno:} $\int \sin(x) \, dx = -\cos(x) + C$
    \item \textbf{Coseno:} $\int \cos(x) \, dx = \sin(x) + C$
\end{itemize}

\section{Tecniche di Integrazione}
Quando la funzione non è immediata, si utilizzano tecniche algebriche per ricondurla a forme note.

\subsection{Integrazione per Parti}
Si basa sulla formula della derivata del prodotto. La regola è:
\begin{equation}
    \int f(x)g'(x) \, dx = f(x)g(x) - \int f'(x)g(x) \, dx
\end{equation}
Questa tecnica è particolarmente utile quando l'integranda è il prodotto di funzioni di diversa natura (es. $x e^x$ o $x \ln(x)$).

\subsection{Integrazione per Sostituzione}
Questa tecnica permette di cambiare la variabile di integrazione per semplificare l'espressione:
\begin{equation}
    \int f(g(x))g'(x) \, dx = \int f(u) \, du \quad \text{con } u = g(x)
\end{equation}
